% ===== PRÉAMBULE CANONIQUE — à coller VERBATIM en tête de CHAQUE .tex =====
\documentclass[11pt,a4paper]{article}
\usepackage[utf8]{inputenc}\usepackage[T1]{fontenc}\usepackage{lmodern}
\usepackage[french]{babel}
\usepackage{amsmath,amssymb,mathtools}\usepackage{icomma}
\usepackage[version=4]{mhchem}          % \ce{...} : équations & formules
\usepackage{siunitx}\sisetup{locale=FR} % \qty{}{}, \si{}, \ang{}
\usepackage{chemfig}                    % \chemfig{...} : structures développées
\usepackage{xcolor}\usepackage{colortbl}
\usepackage{tikz}\usepackage{pgfplots}\pgfplotsset{compat=1.18}
\usepackage[most]{tcolorbox}\usepackage{enumitem}\usepackage{array,booktabs}
\usepackage[a4paper,margin=2.2cm]{geometry}
\usetikzlibrary{arrows.meta,calc,angles,quotes,positioning,patterns,backgrounds,shapes.geometric}
\definecolor{primary}{RGB}{222,49,99}\definecolor{primarydark}{RGB}{150,22,55}
\definecolor{defcol}{RGB}{0,114,178}\definecolor{methcol}{RGB}{52,92,148}
\definecolor{excol}{RGB}{0,135,90}\definecolor{attncol}{RGB}{200,40,55}
\tcbset{boxbase/.style={enhanced,breakable,boxrule=0.9pt,arc=2.5pt,
  left=3mm,right=3mm,top=2.3mm,bottom=2.3mm,fonttitle=\bfseries,coltitle=white}}
% --- boîtes TUTORIEL (noms EXACTS, ne pas renommer) ---
\newtcolorbox{cle}[1][]{boxbase,colback=defcol!6,colframe=defcol,
  title={\textsc{À retenir}\ifx\relax#1\relax\relax\else\textmd{ : \itshape #1}\fi}}
\newtcolorbox{methode}[1][]{boxbase,colback=methcol!7,colframe=methcol,sharp corners,
  title={\textsc{Méthode}\ifx\relax#1\relax\relax\else\textmd{ --- \itshape #1}\fi}}
\newtcolorbox{exemplebox}[1][]{boxbase,colback=excol!5,colframe=excol,
  title={\textsc{Exemple}\ifx\relax#1\relax\relax\else\textmd{ : \itshape #1}\fi}}
\newtcolorbox{piege}[1][]{boxbase,colback=attncol!6,colframe=attncol,
  title={\textsc{Piège}\ifx\relax#1\relax\relax\else\textmd{ : \itshape #1}\fi}}
\newtcolorbox{remarque}[1][]{boxbase,colback=black!4,colframe=black!45,
  title={\textsc{Remarque}\ifx\relax#1\relax\relax\else\textmd{ : \itshape #1}\fi}}
% --- boîtes EXERCICE / CORRIGÉ (noms EXACTS, auto-comptées) ---
\newtcolorbox[auto counter]{exo}[1][]{boxbase,colback=primary!4,colframe=primary,
  title={\textsc{Exercice}~\thetcbcounter\ifx\relax#1\relax\relax\else\textmd{ --- \itshape #1}\fi}}
\newtcolorbox[auto counter]{corrige}[1][]{boxbase,colback=excol!4,colframe=excol,
  title={\textsc{Corrigé}~\thetcbcounter\ifx\relax#1\relax\relax\else\textmd{ --- \itshape #1}\fi}}
\newcommand{\R}{\mathbb{R}}\newcommand{\N}{\mathbb{N}}\newcommand{\Z}{\mathbb{Z}}\newcommand{\ds}{\displaystyle}
% (un \usepackage{fancyhdr}+en-tête est possible ; il est ignoré côté web)
% =========================================================================
\begin{document}
\begin{center}\color{primarydark}\large\bfseries Thème 3 — Corrigé \quad$\bullet$\quad Niveau Facile\end{center}

\begin{corrige}[classer les oxydants]
Le pouvoir oxydant croît avec $E_{\mathrm{r}}^{\circ}$. Du plus faible au plus fort :
\[
\ce{Zn^{2+}}\,(-0{,}76) < \ce{Cu^{2+}}\,(+0{,}34) < \ce{Ag+}\,(+0{,}80) < \ce{F2}\,(+2{,}87).
\]
\ce{F2} est l'oxydant le plus puissant.
\end{corrige}

\begin{corrige}[le plus fort oxydant, le plus fort réducteur]
\begin{enumerate}[label=\alph*)]
  \item L'oxydant le plus fort est celui du couple de $E_{\mathrm{r}}^{\circ}$ le plus \textbf{élevé} :
        \ce{Cl2} ($+1{,}36$~V).
  \item Le réducteur le plus fort est la forme réduite du couple de $E_{\mathrm{r}}^{\circ}$ le plus
        \textbf{bas} : \ce{Al} ($-1{,}66$~V).
\end{enumerate}
\end{corrige}

\begin{corrige}[réaction spontanée ou non ?]
$\Delta E^{\circ}=E_{\mathrm{r}}^{\circ}(\text{oxydant})-E_{\mathrm{r}}^{\circ}(\text{réducteur})$.
\begin{enumerate}[label=\alph*)]
  \item \ce{Cu^{2+}} (oxydant, $+0{,}34$) oxyde \ce{Zn} (réducteur, $-0{,}76$) :
        $\Delta E^{\circ}=0{,}34-(-0{,}76)=+1{,}10>0$. Réaction \textbf{spontanée}.
  \item C'est la réaction inverse : $\Delta E^{\circ}=-1{,}10<0$. \textbf{Non spontanée}
        (\ce{Zn^{2+}} n'oxyde pas \ce{Cu}).
\end{enumerate}
\end{corrige}

\begin{corrige}[quel métal réagit avec un acide ?]
\ce{H+} oxyde le métal (dégagement de \ce{H2}) seulement si
$E_{\mathrm{r}}^{\circ}(\ce{M^{n+}}/\ce{M})<0$.
\begin{enumerate}[label=\alph*)]
  \item \ce{Zn} : $-0{,}76<0$ $\Rightarrow$ \textbf{oui} : \ce{Zn + 2 H+ -> Zn^{2+} + H2 ^}.
  \item \ce{Fe} : $-0{,}44<0$ $\Rightarrow$ \textbf{oui} : \ce{Fe + 2 H+ -> Fe^{2+} + H2 ^}.
  \item \ce{Cu} : $+0{,}34>0$ $\Rightarrow$ \textbf{non} (le cuivre n'est pas attaqué par \ce{HCl}).
  \item \ce{Ag} : $+0{,}80>0$ $\Rightarrow$ \textbf{non}.
\end{enumerate}
\end{corrige}

\begin{corrige}[écrire la réaction spontanée]
\begin{enumerate}[label=\alph*)]
  \item Le couple \ce{Ag+}/\ce{Ag} est le plus haut ($+0{,}80>+0{,}34$) : \ce{Ag+} est l'oxydant.
  \item \ce{Ag+} oxyde \ce{Cu} ($\Delta E^{\circ}=0{,}80-0{,}34=+0{,}46>0$) :
        \[\ce{2 Ag+ + Cu -> 2 Ag + Cu^{2+}}\]
\end{enumerate}
\end{corrige}

\end{document}
